\scanpage{89}
\mspage{44}
the last line in (26) is also identified with
\[
(27)\qquad 1\unsure{-|\longrightarrow}\pi\unsure{-|\longrightarrow}E_{U_s}\longrightarrow E_{k(s)}\unsure{-|\longrightarrow}1
\]
associated with $U_s$, fiber of $U/S$ at $s$, over the residue field $k(s)$.

\uline{Let us summarize} The situation of the relative curve $U/S$, and of an $s\in S$ of codimension $1$, \unsure{defines|defines} by henselization (\inserted{\uline{not}} strict) at $s$ the diagram (2\deleted{\unsure{6|5}}\inserted{\unsure{5|6}}), whose last line is canonically isomorphic to (27) (by \deleted{\unsure{the inclusion}}\inserted{the homomorphism} \inserted{on \unsure{the|the} $\pi_1$} corresponding :
\diagram{\begin{tikzcd}[ampersand replacement=\&]
U_s \arrow[r,no head] \& U_{\widetilde O} \arrow[dd] \\
\unsure{\vert|\downarrow} \& \\
\unsure{s|k} \arrow[r] \& \operatorname{Spec}\widetilde O
\end{tikzcd}}
)

A rational section \inserted{$f_{\deleted{\unsure{K}}}$} of $U$ over $S$, more generally \inserted{$\widetilde f$} of $U_{\widetilde O}$ over $\widetilde S$, defines a section \inserted{$E_{\widetilde f}$} of the middle line of (26), or what amounts to the same a lift
\diagram{\begin{tikzcd}[ampersand replacement=\&,column sep=large]
(28) \& E_{\widetilde K} \arrow[r,"\lambda"] \& E_{U_{\widetilde O}} \\
 \& \& E_{U_s} \arrow[u,"\sim"']
\end{tikzcd}}
of $E_{\widetilde K}\longrightarrow E_{\widetilde O}$ ($\xleftarrow{\sim}E_{k(s)}$) --- where $E_{\widetilde K}$ \unsure{can|by} moreover \unsure{be identified|is identified} with the decomposition group of $\unsure{\underline O_{S,s}|\underline U_{S,s}}$ in $\overline K$. The following conditions are
\margin{88}
